% ========================================================
% BÀI TẬP TOÁN 5 - KẾT NỐI TRI THỨC VỚI CUỘC SỐNG
% Bài 18: Luyện tập chung
% Định dạng: BTTeX Part III - Trắc nghiệm điền số ngắn
% ========================================================

% Câu 1
\begin{ex}
  Đổi $4\text{ km}^2$ thành bao nhiêu héc-ta?
  \shortans{400}
  \loigiai{
    $4\text{ km}^2 = 400\text{ ha}$.
  }
\end{ex}

% Câu 2
\begin{ex}
  Đổi $25\text{ ha}$ thành bao nhiêu mét vuông? (Điền số nghìn mét vuông).
  \shortans{250}
  \loigiai{
    $25\text{ ha} = 250\,000\text{ m}^2 = 250$ nghìn mét vuông.
  }
\end{ex}

% Câu 3
\begin{ex}
  Đổi $8\text{ m}^2\text{ } 50\text{ dm}^2$ thành bao nhiêu đề-xi-mét vuông?
  \shortans{850}
  \loigiai{
    $8\text{ m}^2\text{ } 50\text{ dm}^2 = 850\text{ dm}^2$.
  }
\end{ex}

% Câu 4
\begin{ex}
  Đổi $15\text{ dm}^2\text{ } 4\text{ cm}^2$ thành bao nhiêu xăng-ti-mét vuông?
  \shortans{1504}
  \loigiai{
    $15\text{ dm}^2 = 1\,500\text{ cm}^2$. Do đó kết quả là $1\,504\text{ cm}^2$.
  }
\end{ex}

% Câu 5
\begin{ex}
  Đổi $600\text{ ha}$ thành bao nhiêu ki-lô-mét vuông?
  \shortans{6}
  \loigiai{
    $600 : 100 = 6\text{ km}^2$.
  }
\end{ex}

% Câu 6
\begin{ex}
  Một thửa ruộng hình chữ nhật có chiều dài $80\text{ m}$, chiều rộng $50\text{ m}$. Diện tích thửa ruộng là bao nhiêu a (biết $1\text{ a} = 100\text{ m}^2$)?
  \shortans{40}
  \loigiai{
    Diện tích là $80 \times 50 = 4\,000\text{ m}^2 = 40\text{ a}$.
  }
\end{ex}

% Câu 7
\begin{ex}
  Một khu đất hình chữ nhật có chiều dài $250\text{ m}$, chiều rộng $200\text{ m}$. Diện tích khu đất đó là bao nhiêu héc-ta?
  \shortans{5}
  \loigiai{
    Diện tích là $250 \times 200 = 50\,000\text{ m}^2 = 5\text{ ha}$.
  }
\end{ex}

% Câu 8
\begin{ex}
  Người ta lát sàn một hội trường hình chữ nhật có chiều dài $20\text{ m}$, chiều rộng $12\text{ m}$ bằng các viên gạch vuông cạnh $40\text{ cm}$. Cần dùng tất cả bao nhiêu viên gạch?
  \shortans{1500}
  \loigiai{
    Diện tích hội trường là $20 \times 12 = 240\text{ m}^2 = 2\,400\,000\text{ cm}^2$. Diện tích 1 viên gạch là $40 \times 40 = 1\,600\text{ cm}^2$. Số gạch cần là: $2\,400\,000 : 1\,600 = 1\,500$ (viên).
  }
\end{ex}

% Câu 9
\begin{ex}
  Một khu rừng bảo tồn có diện tích $12\text{ km}^2$. Người ta khoanh vùng $\frac{1}{3}$ diện tích để làm khu bảo tồn động vật hoang dã. Diện tích khu bảo tồn đó là bao nhiêu héc-ta?
  \shortans{400}
  \loigiai{
    $12\text{ km}^2 = 1\,200\text{ ha}$. Diện tích khu bảo tồn là $1\,200 \times \frac{1}{3} = 400\text{ ha}$.
  }
\end{ex}

% Câu 10
\begin{ex}
  Một mảnh đất hình vuông có chu vi là $200\text{ m}$. Diện tích mảnh đất đó là bao nhiêu a (biết $1\text{ a} = 100\text{ m}^2$)?
  \shortans{25}
  \loigiai{
    Cạnh hình vuông là $200 : 4 = 50\text{ m}$. Diện tích là $50 \times 50 = 2\,500\text{ m}^2 = 25\text{ a}$.
  }
\end{ex}

% Câu 11
\begin{ex}
  Đổi $3{,}5\text{ ha}$ thành bao nhiêu mét vuông?
  \shortans{35000}
  \loigiai{
    $3{,}5 \times 10\,000 = 35\,000\text{ m}^2$.
  }
\end{ex}

% Câu 12
\begin{ex}
  Đổi $0{,}04\text{ km}^2$ thành bao nhiêu héc-ta?
  \shortans{4}
  \loigiai{
    $0{,}04 \times 100 = 4\text{ ha}$.
  }
\end{ex}

% Câu 13
\begin{ex}
  Một cánh đồng hình chữ nhật có chu vi $1\,800\text{ m}$, chiều rộng bằng $\frac{4}{5}$ chiều dài. Diện tích cánh đồng là bao nhiêu héc-ta?
  \shortans{20}
  \loigiai{
    Nửa chu vi là $900\text{ m}$. Chiều dài là $900 : (4 + 5) \times 5 = 500\text{ m}$. Chiều rộng là $400\text{ m}$. Diện tích là $500 \times 400 = 200\,000\text{ m}^2 = 20\text{ ha}$.
  }
\end{ex}

% Câu 14
\begin{ex}
  Trên một thửa ruộng rộng $3\text{ ha}$, người ta thu hoạch được $18\text{ tấn}$ thóc. Hỏi trung bình mỗi héc-ta thu hoạch được bao nhiêu tạ thóc?
  \shortans{60}
  \loigiai{
    $18\text{ tấn} = 180\text{ tạ}$. Mỗi héc-ta thu được: $180 : 3 = 60\text{ tạ}$.
  }
\end{ex}

% Câu 15
\begin{ex}
  Một bức tranh hình vuông có chu vi $16\text{ dm}$. Diện tích bức tranh là bao nhiêu xăng-ti-mét vuông?
  \shortans{1600}
  \loigiai{
    Cạnh là $16 : 4 = 4\text{ dm} = 40\text{ cm}$. Diện tích là $40 \times 40 = 1\,600\text{ cm}^2$.
  }
\end{ex}

% Câu 16
\begin{ex}
  Một tấm bìa hình bình hành có đáy $25\text{ cm}$, chiều cao $12\text{ cm}$. Diện tích tấm bìa là bao nhiêu xăng-ti-mét vuông?
  \shortans{300}
  \loigiai{
    Diện tích là $25 \times 12 = 300\text{ cm}^2$.
  }
\end{ex}

% Câu 17
\begin{ex}
  Một miếng tôn hình thoi có hai đường chéo dài $6\text{ dm}$ và $4\text{ dm}$. Diện tích miếng tôn là bao nhiêu đề-xi-mét vuông?
  \shortans{12}
  \loigiai{
    Diện tích là $(6 \times 4) : 2 = 12\text{ dm}^2$.
  }
\end{ex}

% Câu 18
\begin{ex}
  Một công viên hình chữ nhật có chiều dài $1\text{ km}$, chiều rộng $400\text{ m}$. Diện tích công viên là bao nhiêu héc-ta?
  \shortans{40}
  \loigiai{
    $1\text{ km} = 1\,000\text{ m}$. Diện tích là $1\,000 \times 400 = 400\,000\text{ m}^2 = 40\text{ ha}$.
  }
\end{ex}

% Câu 19
\begin{ex}
  Đổi $5\text{ m}^2\text{ } 4\text{ cm}^2$ thành bao nhiêu xăng-ti-mét vuông?
  \shortans{50004}
  \loigiai{
    $5\text{ m}^2 = 50\,000\text{ cm}^2$. Kết quả là $50\,004\text{ cm}^2$.
  }
\end{ex}

% Câu 20
\begin{ex}
  Một sân chơi hình chữ nhật có diện tích $360\text{ m}^2$, chiều rộng $12\text{ m}$. Chu vi sân chơi đó là bao nhiêu mét?
  \shortans{84}
  \loigiai{
    Chiều dài là $360 : 12 = 30\text{ m}$. Chu vi là $(30 + 12) \times 2 = 84\text{ m}$.
  }
\end{ex}

% Câu 21
\begin{ex}
  Một thửa đất hình thang có đáy lớn $40\text{ m}$, đáy bé $20\text{ m}$ và chiều cao $15\text{ m}$. Diện tích thửa đất là bao nhiêu mét vuông?
  \shortans{450}
  \loigiai{
    Diện tích hình thang là: $(40 + 20) \times 15 : 2 = 60 \times 15 : 2 = 450\text{ m}^2$.
  }
\end{ex}

% Câu 22
\begin{ex}
  Một khu đô thị có diện tích $80\text{ ha}$. Diện tích cây xanh và hồ nước chiếm $\frac{3}{10}$ diện tích. Diện tích đất cây xanh và hồ nước là bao nhiêu héc-ta?
  \shortans{24}
  \loigiai{
    Diện tích là $80 \times \frac{3}{10} = 24\text{ ha}$.
  }
\end{ex}

% Câu 23
\begin{ex}
  Đổi $75\,000\text{ dm}^2$ thành bao nhiêu mét vuông?
  \shortans{750}
  \loigiai{
    $75\,000 : 100 = 750\text{ m}^2$.
  }
\end{ex}

% Câu 24
\begin{ex}
  Một kho lương thực có sàn hình chữ nhật diện tích $400\text{ m}^2$. Người ta kê các pallet chứa thóc, mỗi pallet chiếm diện tích $2\text{ m}^2$. Có thể kê tối đa bao nhiêu pallet nếu diện tích lối đi chiếm $60\text{ m}^2$?
  \shortans{170}
  \loigiai{
    Diện tích kê pallet là $400 - 60 = 340\text{ m}^2$. Số pallet tối đa là $340 : 2 = 170$ (pallet).
  }
\end{ex}

% Câu 25
\begin{ex}
  Nếu tăng cả chiều dài và chiều rộng của một hình chữ nhật lên gấp đôi thì diện tích của nó tăng lên gấp mấy lần?
  \shortans{4}
  \loigiai{
    Diện tích mới bằng $2 \times 2 = 4$ lần diện tích ban đầu.
  }
\end{ex}
